\chapter{实数.}

对大小的每一次量度都蕴含着实数的一个弥散的观念。从数学的观点看，实数理论的起源必须追溯到巴比伦科学中一种记数系统逐步形成的过程，这种系统（原则上）能够把值标写得任意接近每一个实数{[}232{]}。掌握了这样一种系统，以及由此必然随之而来的对数值计算的信心，事实上不可避免地以实数的一种“朴素”概念而告终；这种概念与今天在初等教学中、或在物理学家和工程师那里（与十进记数法相联系）重新见到的概念几乎没有差别；这一概念不容许被精确地定义，但可以这样来表述：一个数被认为是通过获得近似值并把它们纳入计算的可能性来定义的：此外，这必然蕴含着由经验给出的、自然不容许无限逼近的大小的量度，与诸如 \(\sqrt{2}\) 之类的“数”（假定存在无限逼近这个数的算法）之间的某种程度的混淆。

因此，类似的“实用主义”观点会在一切以计算技巧取代对严格性的追求和理论关切的数学学派中重新出现。相反，正是后面这些东西支配着希腊数学：我们得感谢它给出了第一个严格而连贯的量的比的理论，也就是说，本质上是实数的理论；它是一系列关于比、特别是关于不可公度比的发现的最终结果，这些发现在希腊思想史上的重要性怎么强调也不为过，但由于缺乏确切的文本，我们只能辨认出其主要轮廓。希腊数学在其开端与关于比例、相似和比的思辨——部分是科学的，部分是哲学的和神秘的——不可分割地联系在一起，特别是与“简单比”（可用分子分母都很小的分数表示的比）的思辨；而宣称用整数和整数之比解释一切，正是毕达哥拉斯学派的诸种特色倾向之一。但恰恰是毕达哥拉斯学派发现了正方形的边与其对角线的不可公度性（\(\sqrt{2}\) 的无理性）：这无疑是数学中不可能性证明的第一个例子；仅仅陈述这个问题这一事实，就已蕴含着比与它的近似值之间的明确区分，并足以指出把希腊数学家与其前辈分隔开来的巨大鸿沟。\(^{1}\)

关于伴随这一重要发现及其后继的思想运动，我们知之甚少。\(^{2}\)我们只打算概括地指出量的比理论赖以建立的主要思想；这个理论由大数学家欧多克索斯（柏拉图的同时代人和朋友）整理定型，被古典希腊数学最终采纳，并通过欧几里得的《几何原本》{[}107{]}为我们所知，它在其中（《几何原本》第 V 卷）得到了大师式的阐述：

\begin{enumerate}
\def\labelenumi{\arabic{enumi})}
\item
  数这个词和数的观念被严格保留给 \textgreater{} 1 的整数（1 是单子，而不是严格意义上的数），这不仅把我们的无理数排除在外，甚至也把我们称为有理数的那些数排除在外，因为对古典时代的希腊数学家来说，它们只是数之比。这里远不止一个单纯的术语问题：“数”这个词对希腊人（以及对直到近代的现代人）来说，是与具有双重合成法则（加法与乘法）的系统的观念联系在一起的：古典希腊数学家把整数之比设想为算子，定义在整数集或其一部分之上（p 与 q 之比是这样的算子：当 N 是 q 的倍数时，它让 N 对应于整数 \(p.(N/q)\)），这些算子构成一个乘法群，但不是具有双重合成法则的系统。在这一点上，希腊数学家自觉地与“计算术家”即职业计算者划清界限，后者像他们的埃及和巴比伦前辈一样，毫无顾忌地把分数或整数与分数之和当作数来处理。此外，他们之所以把这一限制加于数的观念，似乎更多是出于哲学的而非数学的动机，是承袭最早一批希腊思想家关于一与多的思考的结果，因为在这种思想体系中，单位一经细分，便会因这一行为本身而失去其作为单位的性质。\(^{3}\)
\item
  量的理论以公理化方式建立，而且同时对一切种类的量建立（其中可以找到对先前一些理论的暗示，那些理论看来是分别处理长度、面积、体积、时间等的）。某一类量的特征在于：它们可以比较（也就是说，假定了一个严格说来是等价关系的相等定义，以及关系 \textgreater{} 和 \textless{}）、可以加减（\(A + B\) 有定义，而且当 A \textgreater{} B 时 A - B 也有定义），并满足所谓“阿基米德”公理；这条公理从一开始就被清楚地构想为整个大厦的拱顶石（事实上，对实数的一切公理化刻画，它都是不可缺少的）；它得到阿基米德公理这个名字纯属偶然，而阿基米德本人在其《抛物线求积》的引言（{[}5 b{]}，第 II 卷，第 265 页）中就强调：这条公理已被他的前辈们使用过，它在欧多克索斯的工作中起着本质的作用，而且它的推论并不比那些未借助它而作出的面积和体积测定更缺少确定性。\(^{4}\)
\end{enumerate}

容易看出，实数的理论必然从这个公理基础中得出。应当指出，对欧多克索斯来说，给定种类的量构成一个具有单一内部合成法则（加法）的系统，但这个系统还有一个以量的比为算子的外部合成法则，而这些比被设想为构成一个交换的乘法群。设 A 与 \(A'\) 是同一种类的量，B 与 \(B'\) 亦然；A 与 \(A'\) 之比和 B 与 \(B'\) 之比被定义为相等，如果无论整数 m 与 \(m'\) 为何，\(mA < m'A'\) 都蕴含 \(mB < m'B'\)，且 \(mA > m'A'\) 都蕴含 \(mB > m'B'\)；比之间的不等式以类似的方式定义。这些比对每一种类的量都构成一个算子域，这一点等价于第四比例项存在的公理（欧几里得的行文中未明确陈述但经常使用）：给定了比 \(A/A'\) 并给定了 \(B'\)，就存在一个与 \(B'\) 同种类的 B，使得 \(B/B' = A/A'\)。这样，欧多克索斯那个极具独创性的想法使得对各种类的量分别定义的算子域得以相互等同；⁵类似地，整数集（见前文）可以等同于量的比集的一个子集，即有理比（可公度量之比）的集合；然而，由于这些比作为整数上的算子（一般）只对整数集的一个子集有定义，所以仍有必要对它们的理论单独加以展开（欧几里得第 VII 卷）。

这样构造出来的普遍算子域，对希腊数学家来说就相当于我们所说的实数集；此外同样清楚的是，凭借量的加法和量的比的乘法，他们拥有了相当于我们的实数域的东西，尽管其形式远不如我们的便利。\(^{6}\)另一方面，人们也可以自问，他们是否曾把这些集合（给定种类的量的集合，或量的比的集合）设想为我们意义上的完备集合；否则的话，就不容易明白他们为什么会（甚至不觉得有必要把它立为公理）假定第四比例项的存在；再者，某些文本看来涉及这类想法；最后，他们当然把“一条可由连续运动描画的曲线，不能从直线的一侧过渡到另一侧而不与直线相交”视为显然，这一原则他们曾在关于倍立方的研究（通过曲线的交作出 \(\sqrt[3]{2}\)）中使用过，而且它与所考虑的性质本质上等价；然而，我们掌握的文本不容许我们完全精确地了解他们在这一点上的想法。

这就是希腊数学古典时代实数理论的状况。欧多克索斯的构造尽管令人赞叹，从连贯性和严格性的观点看无可指摘，但必须承认它缺乏柔性，很不利于数值计算尤其是代数计算的发展。更何况，它的逻辑必然性只有崇尚严格、精于抽象的头脑才能领会；因此，随着希腊数学的衰落，人们看到由计算术家的传统保留下来的“朴素”观点逐渐重新露头，这很自然；例如在丢番图{[}91 a{]}那里占上风的就是这种观点，他是这一传统真正的传人，而非官方希腊科学的代表；后者虽然出于形式上的考虑复述了欧几里得的数的定义，但他用“数”这个词真正指的，是代数问题中的未知数，其解或为整数，或为分数，甚至为无理数。\(^{7}\)尽管关于数这一主题的态度上的这种变化，联系着数学史上最重要的进展之一，即代数的发展（见第 48 页），但必须明白，它本身并不构成进步，反而是退步。

我们不可能在这里追述直到中世纪末为止数的观念在印度、阿拉伯和西方数学中经历的一切变迁：占据统治地位的是数的“朴素”概念；而且，尽管这一时期欧几里得的《几何原本》被用作数学教学的基础，欧多克索斯的学说很可能普遍不曾被人理解，因为对它的需要已不再出现。欧几里得的“比”最常被冠以“数”之名；整数的计算法则被应用于它们，并由此得到正确的结果，却不去深入剖析这些方法何以成功的原因。

然而我们看到，早在 16 世纪中叶，R. 邦贝利就在其《代数学》{[}28 b{]}\(^{8}\)中就这一问题写下了一个观点，它（在假定欧几里得第 V 卷的结果已经掌握的条件下）本质上是正确的；他认识到，一旦选定了长度单位，长度与量的比之间便存在一个双射对应，于是他在长度上定义各种代数运算（当然假定单位固定），并且用长度表示数，从而得到了实数域的几何定义（这一观点的功劳通常归于笛卡尔），这样就给他的《代数学》奠定了一个坚实的几何基础。\(^{9}\)

但邦贝利的《代数学》尽管对那个时代而言先进得出奇，也没有超出根式开方和二、三、四次方程的根式求解；当然，开根式的可能性被他毫无讨论地接受了下来。西蒙·斯蒂文{[}295{]}在数的问题上也采取类似的观点：数对他而言就是标记量之量度的东西，而且他认为它本质上是“连续的”（并未明确他对这个词所赋予的含义）；即使他区分“几何数”和“算术数”，那也只是出于其定义方式的偶然，在他看来二者并无种类之别；下面就是他关于这个问题的最后的话：“因此我们断定，不存在什么荒谬的、无理的、不规则的、不可言喻的或聋的数；数之中有如此卓越与和谐，以至它们那令人赞叹的完美给了我们日夜沉思的题材”（{[}295{]}，第 10 页）。另一方面，由于第一个把十进小数这件工具变成一种计算方法，并为它们提出了一种已接近我们今天所用的记号，他清楚地认识到这些小数为每个实数提供了一个无限逼近的算法，正如他在 1594 年的那本题为“包含一条适用于一切方程的通则”的《代数附录》（该小册子唯一已知的一册于 1914 年在鲁汶被焚毁；但可参见{[}295{]}，第 88 页）中所表明的。把这样一个方程写成 \(P(x) = Q(x)\) 的形式（其中 \(P\) 是次数高于多项式 \(Q\) 的多项式，且 \(P(0) < Q(0)\)）后，我们对 \(x\) 代入 10、100、1000、\ldots{} 直到发现 \(P(x) > Q(x)\)，他说，这就确定了根的位数；然后（比如若发现根有两位）代入 10、20、\ldots{} 这就确定了十位上的数字；接着对下一位数字如法炮制，然后对相继的各位小数也是如此：“如此无限进行下去”，他说，“就会无限接近所求者”（{[}295{]}，第 88 页）。可以看出，斯蒂文（无疑是最早）对波尔查诺定理有了清晰的想法，并且认识到这一定理是数值方程系统求解的本质工具；与此同时，在那里还可以认出一种关于数连续统的直观概念，它如此清晰，以至要把它最终精确化已所剩无几。

然而，在随后的两个世纪里，正确方法的最终确立两度被两门理论的发展所延缓，我们无须在此叙述它们的历史：无穷小演算和级数理论。在它们引发的争论中，我们可以辨认出——如同在数学史的一切时期一样——两类人之间永恒的拉锯：一类是埋头向前的研究者，不惜付出一些不安全的代价，并相信日后总会有时间去巩固业已征服的领土；另一类是批判性的头脑，他们（在直觉能力和发明才干上未必输给前者）并不认为把自己的气力花在概念的精确表述和严格证成上是白费。在 17 世纪，争论的主要对象是无穷小的概念，它凭借所能达到的结果而获得事后的辩护，却似乎与阿基米德公理公开对立；我们看到那个时代最开放的头脑最终采取了一个与邦贝利相去不远的观点，其区别主要在于对古人严格方法给予了更多注意；艾萨克·巴罗（牛顿的老师，他自己也在无穷小演算的创立中起过重要作用）在其 1664-65-66 年间于剑桥讲授的《数学讲义》{[}16 a 与 b{]}中对这一观点作了精彩的阐述；他认识到，要让数的课题重新获得那句格言式的“几何确实性”，就必须回到欧多克索斯的理论，于是他极其公正地长篇为后者辩护（据他见证，这理论在他的许多同时代人看来是可理解的），以驳斥那些以晦涩甚至荒谬相责难的人。另一方面，他把数定义为表示量之比的符号，它们可借助算术运算彼此结合，由此得到了实数域，这些说法后来被牛顿在其《算术》中所承袭，而直到戴德金和康托尔，他的后继者们对之未曾改动分毫。

但恰恰是在这个时期前后，级数展开的方法被引入了，它不久就在积习难改的代数学家手中带上了纯粹形式化的性质，并把数学家的注意力从收敛问题上引开，而在实数的领域中，对级数的健全使用必然会提出收敛问题。这一方法的主要创立者牛顿仍然意识到必须考虑这些问题：而且，即使他没有把它们充分阐明，他至少已认识到，就变量的小值而言，他所引入的幂级数“在大多数情况下”至少收敛得同几何级数（其收敛性古人心知肚明）一样好（{[}233 a{]}，第 I 卷，第 3–26 页）；大约在同一时期，莱布尼茨已观察到，各项绝对值递减且趋于 0 的交错级数是收敛的；到下一个世纪，达朗贝尔于 1768 年对非收敛级数的使用表示了怀疑。但伯努利、尤其是欧拉的权威，使这类怀疑在当时只是例外。

显然，习惯于在数值计算中使用级数的数学家决不会这样忽视收敛的概念；而在这门领域中如同在许多其他领域中一样，第一个带来正确方法之回归的人，是一位从青年时代起就热爱数值计算的数学家，这并非出于偶然：C. F. 高斯几乎在孩提时代就使用过算术几何平均的算法，\(^{10}\)因而不能不为自己形成清晰的极限观念；我们看到他在一篇写于 1800 年（但直到我们时代才发表）的残稿（{[}124 a{]}，第 X \(_1\) 卷，第 390 页）中，一方面精确地定义了实数序列的上确界和下确界，另一方面定义了它的上极限和下极限；前者的存在性（对有界序列）似乎被当作显然而接受，后者则被正确地定义为当 \(n\) 趋于 +∞ 时 \(\sup_{p≥0} u_{n+p}\)、\(\inf_{p≥0} u_{n+p}\) 的极限。高斯还在其 1812 年关于超几何级数的论文（{[}124 a{]}，第 III 卷，第 139 页）中给出了关于收敛性的讨论的第一个典范，照他的说法，这讨论是“以全部的严格性进行的，为的是让那些偏爱古代几何学家严格方法的人满意”：诚然，这项讨论在该论文中只构成一个次要之点，并没有追溯到级数理论的第一原理；是柯西在其 1821 年的《分析教程》（{[}56 a{]}，(2)，第 III 卷）中最先建立了这些原理，其方式在一切点上都正确，从明确陈述并被视为显然的柯西准则出发；由于在数的定义上他拘守巴罗和牛顿的观点，因此可以说，对他而言，实数是由量的公理和柯西准则来定义的：这确实足以定义它们。

也是在这个时期，实数理论的另一个重要方面得到了最终的澄清。如前所述，两条连续曲线不能彼此穿越而不相交，这一点一直被当作几何上显然的事而接受；这条原则（经过适当的精确化）也将等价于直线是完全空间这一性质。这条原则仍然属于高斯 1799 年对达朗贝尔定理——即每个实系数多项式都有一个实的或复的根——所作“严格”证明的基础（{[}124 a{]}，第 III 卷，第 1 页）；高斯 1815 年给出的同一定理的证明（{[}124 a{]}，第 III 卷，第 31 页），与拉格朗日先前的一次尝试一样，依据的是一条类似但更简单的原则，即多项式不能变号而不变为零（我们已经见过斯蒂文使用这条原则）。1817 年，波尔查诺从柯西准则出发，对后一条原则给出了一个完全的证明，并把它作为一个关于数值变量的连续数值函数的类似定理的特例得到{[}27 c{]}。他（先于柯西）清楚地陈述了“柯西准则”，并试图用一个论证来证成它，而在缺乏实数的一切算术定义的情况下，这个论证只能是、也只可能是一个循环；但是，一旦这一点得到承认，他的工作就完全是正确的而且相当出色，因为其中不仅包含连续函数的现代定义（此处系首次给出）以及多项式连续性的证明，甚至还包含任意有界实数集的下确界存在的证明（他不谈集合，而是谈实数的性质，这二者是一回事）。另一方面，柯西在其《分析教程》（{[}56 a{]}，(2)，第 III 卷）中也定义了一元或多元数值变量的连续函数，同样证明了一元的连续函数不能变号而不变为零，而且用的论证与西蒙·斯蒂文的相同，一旦连续性有了定义，只要使用柯西准则（或者只要像柯西此时那样，假定称为“区间套”的等价原则——无穷十进小数的收敛当然只是它的一个特例），这论证自然就成为正确的。

一旦达到了这一点，数学家们剩下的工作只是通过纠正若干错误、填补若干空缺，把已获得的结果加以精确化和发展。例如，柯西一度相信，各项都是变量的连续函数的收敛级数，其和也是一个连续函数：阿贝尔在其关于级数的重要工作（{[}1{]}，第 I 卷，第 219 页；另参见第 II 卷，第 257 页及以下各处）中对这一点的纠正，最终以魏尔斯特拉斯在其讲课（未出版，但影响巨大）中对一致收敛概念的阐明而告结束（见第 205 页）。另一方面，柯西在他给出的多项式根存在性的一个证明中，曾没有充分根据地假定了连续函数最小值的存在；又是

魏尔斯特拉斯在其讲课中通过证明定义在闭有界区间上的数值变量函数的最小值存在，使这类问题得到了澄清；正是从他对把这一定理不加证成地应用于函数集（其中“狄利克雷原理”是最著名的例子）的批评出发，开始了那场思想运动，它最终（见第 143 页）以紧空间的一般定义和该定理的现代陈述而告终。

与此同时，魏尔斯特拉斯在其讲课中认识到，把实数的观念与量的理论完全分离开来在逻辑上是有意义的：这样做的确可以归结为以公理化方式定义直线上的点集（从而最终定义实数集）并假定这样一个集合存在；尽管这种做法本质上是正确的，但显然更可取的是只从有理数出发，然后通过完备化从那里导出实数。\(^{11}\)魏尔斯特拉斯、戴德金、梅雷和康托尔以不同的方法、彼此独立地做到了这一点；戴德金提出的所谓“分割”法（{[}79{]}，第 II 卷，第 315–334 页）与欧多克索斯的诸定义接近得多，而其他几种提出的方法则较接近于豪斯多夫以后用来完备化一个度量空间的方法。也正是在这个时刻，康托尔开始发展实数集的理论——戴德金早就有了关于它的最初想法{[}48{]}——并由此得到了关于直线拓扑、开集与闭集的结构、导集、完全不连通的完全集等概念的主要初等结果。

到康托尔为止，实数理论除去少数例外已具有其最终形式；我们来简要地指出它是以何种方式被推广的。除一般拓扑方面的工作（见第 143 页）以及对积分论的应用（见第 221 页）之外，这里主要涉及的是关于直线上点集及实变量的数值函数的结构与分类的研究；它们起源于博雷尔的工作{[}32 a{]}，这些工作主要着眼于测度论，但除其他结果外，最终得出了“博雷尔集”的定义：这些集属于 R 的子集中包含诸区间、且对并、可数交以及 C 运算封闭的最小族。与这些集密切相关的是所谓“贝尔”函数，也就是可以从连续函数出发、经“超穷”重复的极限运算而得到的那些函数；贝尔在一项重要工作中定义了它们，在这项工作中他完全抛弃了测度的观点，以便系统地处理这些问题的定性方面和“拓扑”方面{[}11 a{]}：正是在这个场合，他第一次定义并研究了半连续函数，而且为了刻画连续函数的极限函数，他引入了贫集（用贝尔的术语说，“第一纲”集）这一重要概念（见第 165 页）。至于继贝尔之后的大量工作——它们主要出自俄国尤其是波兰学派——我们在这里只能提请注意它们的存在（见第 166 页）。

